\documentclass[11pt,a4paper]{article}

\usepackage[utf8]{inputenc}
\usepackage[english]{babel}
\usepackage{fontspec}
\setmainfont{Times New Roman}
\setsansfont{Arial}
\setmonofont{Consolas}

\usepackage[a4paper,top=1.5cm,bottom=1.6cm,left=1.7cm,right=1.7cm]{geometry}
\usepackage{amsmath,amssymb}
\usepackage{booktabs}
\usepackage{tabularx}
\usepackage{array}
\usepackage{listings}
\usepackage{xcolor}
\usepackage{fancyhdr}
\usepackage{multicol}
\usepackage{enumitem}
\usepackage{tcolorbox}
\usepackage{lastpage}

% Colors for syntax highlighting and frames
\definecolor{codegray}{rgb}{0.96,0.96,0.96}
\definecolor{keywordblue}{rgb}{0.0,0.2,0.6}
\definecolor{stringgreen}{rgb}{0.0,0.5,0.2}
\definecolor{commentgray}{rgb}{0.4,0.4,0.4}
\definecolor{boxborder}{rgb}{0.7,0.7,0.7}

\lstset{
  language=Python,
  backgroundcolor=\color{codegray},
  basicstyle=\ttfamily\small,
  keywordstyle=\color{keywordblue}\bfseries,
  stringstyle=\color{stringgreen},
  commentstyle=\color{commentgray}\itshape,
  showstringspaces=false,
  breaklines=true,
  frame=single,
  rulecolor=\color{boxborder},
  aboveskip=2.5pt,
  belowskip=2.5pt,
  numbers=none,
  columns=fullflexible
}

% Header & Footer
\pagestyle{fancy}
\fancyhf{}
\setlength{\headheight}{14pt}
\lhead{\small \textbf{Course:} Data Analysis with Python (DSAI1005)}
\rhead{\small \textbf{Midterm Exam --- 60 Minutes}}
\lfoot{\small Faculty of Data Science \& AI --- National Economics University}
\rfoot{\small Page \thepage/6}
\renewcommand{\headrulewidth}{0.5pt}
\renewcommand{\footrulewidth}{0.5pt}

\newcommand{\answerbox}[2]{
  \vspace{1mm}
  \begin{tcolorbox}[
    colback=white,
    colframe=gray!60,
    boxrule=0.6pt,
    arc=2pt,
    height=#1,
    left=3mm, right=3mm, top=2mm, bottom=2mm
  ]
  \small\textit{\color{gray}#2}
  \end{tcolorbox}
  \vspace{1mm}
}

\begin{document}

% =========================================================================
% PAGE 1: HEADER + SCORE TABLE + ANSWER GRID + QUESTIONS 1 TO 3
% =========================================================================
\noindent
\begin{minipage}[t]{0.54\textwidth}
  \textbf{NATIONAL ECONOMICS UNIVERSITY}\\
  \textbf{FACULTY OF DATA SCIENCE \& AI}\\
  \rule{4.5cm}{0.4pt}\\[1.5mm]
  \textbf{MIDTERM EXAMINATION (PAPER-BASED)}\\
  \textbf{Course:} Data Analysis with Python (DSAI1005)\\
  \textbf{Scope:} NumPy (20\%), Pandas (20\%), Viz (20\%), Code (40\%)\\
  \textbf{Difficulty:} 60\% Easy/Basic + 40\% Medium\\
  \textbf{Duration:} 60 minutes --- \textbf{Exam Code:} \textbf{301-EN}\\
  \textit{\small (No personal laptops, textbooks, or unauthorized notes allowed)}
\end{minipage}
\hfill
\begin{minipage}[t]{0.43\textwidth}
  \begin{tabular}{|p{2.1cm}|p{2.1cm}|p{2.1cm}|}
    \hline
    \centering \textbf{Section I} & \centering \textbf{Section II} & \centering \textbf{TOTAL} \tabularnewline
    \hline
    & & \\[0.8cm]
    \hline
    \multicolumn{3}{|l|}{\small \textit{Examiner 1:}} \tabularnewline
    \multicolumn{3}{|l|}{\small \textit{Examiner 2:}} \tabularnewline
    \hline
  \end{tabular}
\end{minipage}

\vspace{2mm}
\noindent
\textbf{Full Name:} \dotfill{} \textbf{Student ID:} \dotfill{} \textbf{Exam Code:} \textbf{301-EN}\\
\textbf{Class / Cohort:} \dotfill{} \textbf{Exam Room:} \dotfill{} \textbf{Seat No:} \dotfill

\vspace{2mm}
\noindent
\textbf{MULTIPLE-CHOICE ANSWER SHEET} \textit{(Write letters A, B, C, or D in the corresponding boxes)}:
\begin{center}
\small
\begin{tabular}{|c|c|c|c|c|c|c|c|c|c|c|c|c|}
  \hline
  \textbf{Question} & \textbf{1} & \textbf{2} & \textbf{3} & \textbf{4} & \textbf{5} & \textbf{6} & \textbf{7} & \textbf{8} & \textbf{9} & \textbf{10} & \textbf{11} & \textbf{12} \\
  \hline
  \textbf{Answer} & & & & & & & & & & & & \\[3.5mm]
  \hline
\end{tabular}
\end{center}

\vspace{-3mm}
\hrule
\vspace{1.5mm}

\section*{SECTION I: MULTIPLE-CHOICE QUESTIONS (6.0 pts --- 0.5 pt each)}
\textit{\small (60\% Easy \& Basic Questions: Foundational Knowledge \& Core Syntax)}

\begin{enumerate}[leftmargin=*,itemsep=2.5pt]

\item What array is generated by the NumPy expression \texttt{np.arange(2, 11, 2)}?
\begin{multicols}{4}
\begin{enumerate}[label=\Alph*.]
  \item \texttt{[2, 4, 6, 8]}
  \item \texttt{[2, 4, 6, 8, 10]}
  \item \texttt{[2, 5, 8, 11]}
  \item \texttt{[4, 6, 8, 10]}
\end{enumerate}
\end{multicols}

\item For a two-dimensional array \texttt{a = np.array([[10, 20, 30], [40, 50, 60]])}, what are its shape (\texttt{a.shape}) and number of dimensions (\texttt{a.ndim})?
\begin{multicols}{2}
\begin{enumerate}[label=\Alph*.]
  \item \texttt{shape: (2, 3), ndim: 2}
  \item \texttt{shape: (3, 2), ndim: 2}
  \item \texttt{shape: (6,), ndim: 1}
  \item \texttt{shape: (2, 3), ndim: 6}
\end{enumerate}
\end{multicols}

\item Consider the following 2D array:
\begin{lstlisting}
m = np.array([[1, 2, 3],
              [4, 5, 6],
              [7, 8, 9]])
\end{lstlisting}
Which expression correctly accesses the element with value \texttt{6}?
\begin{multicols}{4}
\begin{enumerate}[label=\Alph*.]
  \item \texttt{m[2, 1]}
  \item \texttt{m[1, 1]}
  \item \texttt{m[1, 2]}
  \item \texttt{m[2, 2]}
\end{enumerate}
\end{multicols}

\end{enumerate}

\newpage
% =========================================================================
% PAGE 2: QUESTIONS 4 TO 8 (NUMPY Q4 + PANDAS Q5 TO Q8)
% =========================================================================
\begin{enumerate}[leftmargin=*,itemsep=2.5pt,resume]

\item For a two-dimensional NumPy array \texttt{arr} of shape \texttt{(4, 5)}, what does the function call \texttt{np.mean(arr, axis=0)} compute?
\begin{enumerate}[label=\Alph*.]
  \item Computes a single overall mean across all elements in the array.
  \item Computes the mean of each column, returning an array of 5 values.
  \item Computes the mean of each row, returning an array of 4 values.
  \item Transposes the array and calculates standard deviations.
\end{enumerate}

\item What is the fundamental structural difference between a Pandas \texttt{Series} and a \texttt{DataFrame}?
\begin{enumerate}[label=\Alph*.]
  \item A \texttt{Series} can store only text, whereas a \texttt{DataFrame} can store only floating-point numbers.
  \item A \texttt{Series} does not support indexing, whereas a \texttt{DataFrame} requires numeric indexes.
  \item There is no difference; the two names are interchangeable synonyms in Pandas.
  \item A \texttt{Series} is a 1-dimensional labeled array, whereas a \texttt{DataFrame} is a 2-dimensional labeled tabular structure of rows and columns.
\end{enumerate}

\item Given DataFrame \texttt{df}, which indexing attribute is specifically designed to select rows by their \textbf{0-based integer position} (e.g., the very first row at index 0)?
\begin{multicols}{4}
\begin{enumerate}[label=\Alph*.]
  \item \texttt{df.iloc[0]}
  \item \texttt{df.loc[0]}
  \item \texttt{df.get[0]}
  \item \texttt{df.row[0]}
\end{enumerate}
\end{multicols}

\item After loading a tabular dataset into DataFrame \texttt{df}, which concise statement is most directly used to count the number of missing values (\texttt{NaN}) in each column?
\begin{multicols}{4}
\begin{enumerate}[label=\Alph*.]
  \item \texttt{df.describe()}
  \item \texttt{df.info(count=0)}
  \item \texttt{df.isna().sum()}
  \item \texttt{df.drop\_na()}
\end{enumerate}
\end{multicols}

\item Which Pandas expression correctly filters rows where column \texttt{'City'} equals \texttt{'Hanoi'} \textbf{and} column \texttt{'Age'} is greater than or equal to \texttt{20}?
\begin{enumerate}[label=\Alph*.]
  \item \texttt{df[(df['City'] == 'Hanoi') \& (df['Age'] >= 20)]}
  \item \texttt{df[(df['City'] == 'Hanoi') and (df['Age'] >= 20)]}
  \item \texttt{df[df['City'] == 'Hanoi' \& df['Age'] >= 20]}
  \item \texttt{df.filter(City == 'Hanoi', Age >= 20)}
\end{enumerate}

\end{enumerate}

\newpage
% =========================================================================
% PAGE 3: QUESTIONS 9 TO 12 (VISUALIZATION)
% =========================================================================
\begin{enumerate}[leftmargin=*,itemsep=2.5pt,resume]

\item In exploratory data analysis, which type of chart is most suitable for visualizing the historical trend of a continuous metric (such as monthly revenue) over time?
\begin{multicols}{4}
\begin{enumerate}[label=\Alph*.]
  \item Pie chart
  \item Line chart
  \item Heatmap
  \item Scatter plot
\end{enumerate}
\end{multicols}

\item Which visualization tool displays the 5-number summary (minimum, first quartile $Q_1$, median, third quartile $Q_3$, maximum) and visually identifies extreme outlier data points?
\begin{multicols}{4}
\begin{enumerate}[label=\Alph*.]
  \item Bar plot (\texttt{sns.barplot})
  \item Count plot (\texttt{sns.countplot})
  \item Pie chart (\texttt{plt.pie})
  \item Box plot (\texttt{sns.boxplot})
\end{enumerate}
\end{multicols}

\item In Seaborn, consider the following relational plotting statement:
\begin{lstlisting}
import seaborn as sns
sns.scatterplot(data=df, x='Income', y='Spending', hue='Gender')
\end{lstlisting}
What is the primary visual role of the \texttt{hue='Gender'} parameter?
\begin{enumerate}[label=\Alph*.]
  \item It splits the figure into two completely separate subplots.
  \item It modifies the opacity (alpha) of the plotted points.
  \item It colors data points differently based on the categories of \texttt{Gender}.
  \item It changes the size of the font used in the axis labels.
\end{enumerate}

\item In Matplotlib, which utility function is commonly called before displaying or saving a figure to automatically adjust subplot spacing and prevent titles or tick labels from overlapping?
\begin{multicols}{4}
\begin{enumerate}[label=\Alph*.]
  \item \texttt{plt.tight\_layout()}
  \item \texttt{plt.autospace()}
  \item \texttt{plt.clean\_plot()}
  \item \texttt{plt.resize()}
\end{enumerate}
\end{multicols}

\end{enumerate}

\newpage
% =========================================================================
% PAGE 4: SECTION II - QUESTION 1 (NUMPY & PANDAS PIPELINE - 2.0 PTS)
% =========================================================================
\section*{SECTION II: APPLIED CODING \& SHORT-ANSWER (4.0 pts)}
\textit{\small (40\% Medium-Level Questions: Integrated Data Processing \& Visual Analytics)}

\subsection*{Question 1: Retail Sales Analysis with NumPy \& Pandas (2.0 pts) --- [Medium]}

\noindent
\textbf{Scenario:} An e-commerce business tracks retail orders. You are given a DataFrame \texttt{df} containing 5 customer transactions:

\begin{lstlisting}
import numpy as np
import pandas as pd

df = pd.DataFrame({
    'OrderID':  ['O01', 'O02', 'O03', 'O04', 'O05'],
    'Category': ['Tech', 'Book', 'Tech', 'Home', 'Book'],
    'Price':    [500.0, 80.0, 1200.0, 300.0, 60.0],
    'Quantity': [2, 5, 1, 3, 10],
    'Discount': [0.10, 0.00, 0.05, 0.10, 0.00]
})
\end{lstlisting}

\noindent
\textbf{Tasks (write exact and concise Python code):}
\begin{enumerate}[label=\textbf{\alph*)}]
  \item \textbf{(0.4 pts)} Write a vectorized Pandas expression to calculate total net revenue for each order and store it in a new column \texttt{df['Revenue']}:
  $$\text{Revenue} = \text{Price} \times \text{Quantity} \times (1 - \text{Discount})$$
  \item \textbf{(0.4 pts)} Using \texttt{np.where}, create a new column \texttt{df['Priority']} that assigns \texttt{'High'} if \texttt{Revenue >= 1000} \textbf{or} \texttt{Quantity >= 5}, otherwise assigns \texttt{'Standard'}.
  \item \textbf{(0.4 pts)} Write code to filter all orders where \texttt{Category} is \texttt{'Tech'} \textbf{and} \texttt{Quantity >= 2}, storing the result in \texttt{tech\_orders}.
  \item \textbf{(0.4 pts)} Group the DataFrame by \texttt{'Category'} and compute both the \textbf{sum} and \textbf{mean} of \texttt{'Revenue'}, storing the result in \texttt{summary}.
  \item \textbf{(0.4 pts)} Sort the resulting summary table \texttt{summary} by total revenue (\texttt{'sum'}) in \textbf{descending} order.
\end{enumerate}

\answerbox{7.2cm}{
\textbf{Work Area --- Question 1:}\\[2mm]
\textbf{a) Compute Revenue column (0.4 pts):}\\
df['Revenue'] = \dotfill\\[4.5mm]
\textbf{b) Conditional Priority classification with np.where (0.4 pts):}\\
df['Priority'] = np.where(\dotfill)\\[4.5mm]
\textbf{c) Filter Tech orders with Quantity >= 2 (0.4 pts):}\\
tech\_orders = df[\dotfill]\\[4.5mm]
\textbf{d) Groupby Category and calculate sum and mean (0.4 pts):}\\
summary = df.groupby(\dotfill)\\[4.5mm]
\textbf{e) Sort summary table by total revenue descending (0.4 pts):}\\
final\_sorted = summary.sort\_values(\dotfill)
}

\newpage
% =========================================================================
% PAGE 5: SECTION II - QUESTION 2 (VISUALIZATION WITH SEABORN - 1.5 PTS)
% =========================================================================
\subsection*{Question 2: Sales Performance Visualization with Matplotlib \& Seaborn (1.5 pts) --- [Medium]}

\noindent
\textbf{Scenario:} Continuing with the retail sales dataset \texttt{df} from Question 1 (with computed column \texttt{'Revenue'}), you need to create a 2-panel figure comparing sales performance across categories.

\begin{lstlisting}
import matplotlib.pyplot as plt
import seaborn as sns
\end{lstlisting}

\noindent
\textbf{Tasks (write correct and standard plotting code):}
\begin{enumerate}[label=\textbf{\alph*)}]
  \item \textbf{(0.5 pts)} Write the Matplotlib Object-Oriented command to create a figure with a $1 \times 2$ subplot grid (1 row, 2 columns) of size $12 \times 4$ inches, storing the results in \texttt{fig} and \texttt{axes}.
  \item \textbf{(1.0 pts)} Complete the code to plot:
  \begin{itemize}
    \item On \texttt{axes[0]}: Render a Seaborn bar plot (\texttt{sns.barplot}) showing average \texttt{'Revenue'} across \texttt{'Category'}.
    \item On \texttt{axes[1]}: Render a Seaborn box plot (\texttt{sns.boxplot}) of \texttt{'Revenue'} across \texttt{'Category'} to inspect data distribution and identify any outliers.
  \end{itemize}
\end{enumerate}

\answerbox{13.0cm}{
\textbf{Work Area --- Question 2:}\\[3mm]
\textbf{a) Canvas initialization with 1 row and 2 columns (0.5 pts):}\\
fig, axes = \dotfill\\[6mm]
\textbf{b) Plotting on subplots (1.0 pts):}\\[2mm]
\textit{\# Subplot 1: Bar chart of average Revenue on axes[0]:}\\
sns.barplot(\dotfill\\[4mm]
\hspace*{2.5cm}\dotfill)\\[6mm]
\textit{\# Subplot 2: Box plot of Revenue on axes[1]:}\\
sns.boxplot(\dotfill\\[4mm]
\hspace*{2.5cm}\dotfill)\\[6mm]
\textit{\# Subplot titles (already provided for reference):}\\
axes[0].set\_title("Average Revenue by Category")\\
axes[1].set\_title("Revenue Distribution \& Outliers by Category")
}

\newpage
% =========================================================================
% PAGE 6: SECTION II - QUESTION 3 (LAYOUT & INSIGHTS - 0.5 PTS) + END
% =========================================================================
\subsection*{Question 3: Layout Optimization \& Chart Interpretation (0.5 pts) --- [Medium]}

\noindent
\textbf{Scenario:} After rendering the two subplots in Question 2, you want to ensure the figure is ready for executive presentation and draw an analytical conclusion.

\noindent
\textbf{Tasks:}
\begin{enumerate}[label=\textbf{\alph*)}]
  \item \textbf{(0.25 pts)} Write the single line of Matplotlib code to automatically adjust padding between subplots, ensuring that titles and axis labels do not overlap.
  \item \textbf{(0.25 pts)} In a box plot such as the one on \texttt{axes[1]}, explain briefly how an analyst can visually determine whether a transaction is classified as an outlier based on the box whiskers.
\end{enumerate}

\answerbox{10.0cm}{
\textbf{Work Area --- Question 3:}\\[3mm]
\textbf{a) Command to optimize layout and prevent label overlap (0.25 pts):}\\[2mm]
\dotfill\\[8mm]
\textbf{b) Visual interpretation of outliers in a box plot (0.25 pts):}\\[3mm]
\dotfill\\[4mm]
\dotfill\\[4mm]
\dotfill\\[4mm]
\dotfill
}

\vspace{5mm}
\begin{center}
\textbf{------------------- END OF EXAM -------------------}\\
\textit{\small (Proctors will not answer questions regarding exam content. Hand in all exam sheets upon completion.)}
\end{center}

\end{document}
